\documentclass[lualatex,paper=b5j,fontsize=10pt]{jlreq}

\usepackage[twocolumn=false,footer=false]{surikumi}

\newcommand{\PatternSample}[3]{%
  \par\noindent{\sffamily\bfseries #1}\par\smallskip
  \begingroup
    \MathMagazineSetup{
      pattern=#2,
      series-position=#3,
      series={ベーシック演習},
      title={数列の基礎},
      author={数学編集部},
      author-font=mincho,
      author-position=right,
      title-fill=black
    }%
    \MMTitleBanner
  \endgroup
  \par\vspace{4mm}%
}

% Variant sampler for the mathematical-motif patterns, which also accepts
% pattern-tone / pattern-accent / pattern-scale settings.
\newcommand{\PatternTuned}[3][]{%
  \par\noindent{\sffamily\bfseries #2}\par\smallskip
  \begingroup
    \MathMagazineSetup{
      pattern=#3,
      series-position=left,
      series={講義／数学},
      title={数列の基礎},
      author={数学編集部},
      author-font=mincho,
      author-position=right,
      title-fill=black,
      #1
    }%
    \MMTitleBanner
  \endgroup
  \par\vspace{4mm}%
}

\newcommand{\PatternVariant}[6]{%
  \par\noindent{\sffamily\bfseries #1}\par\smallskip
  \begingroup
    \MathMagazineSetup{
      pattern=#2,
      series-position=#3,
      series={講義／数 IIAIIB},
      title={集合と論証},
      author={数学編集部},
      author-font=#4,
      author-position=#5,
      title-fill=#6
    }%
    \MMTitleBanner
  \endgroup
  \par\vspace{4mm}%
}

\begin{document}

\begin{center}
  {\sffamily\gtfamily\bfseries\LARGE 背景パターン見本}\par
  \smallskip
  {\small 31 patterns: scales / triangles / diamonds / chevrons / weave / \ldots\\
   和柄 10 種: asanoha / seigaiha / shippou / kikkou / sayagata / \ldots\\
   数学的モチーフ 9 種: truchet / isocubes / sierpinski / lattice / \ldots}
\end{center}
\vspace{3mm}

\PatternSample{1. scales（既定・鱗状）}{scales}{left}
\PatternSample{2. triangles（三角形）}{triangles}{center}
\PatternSample{3. diamonds（菱形）}{diamonds}{left}
\PatternSample{4. chevrons（山形）}{chevrons}{center}

\clearpage

\PatternSample{5. weave（織り目）}{weave}{left}
\PatternSample{6. hexagons（六角形）}{hexagons}{center}
\PatternSample{7. pinwheels（風車）}{pinwheels}{left}
\PatternSample{8. zigzags（ジグザグ）}{zigzags}{center}

\clearpage

\PatternSample{9. checker（市松）}{checker}{left}
\PatternVariant{10. parquet（寄木線画）}{parquet}{center}{gothic}{center}{black}
\PatternVariant{11. origami（折紙線画・白抜き題字）}{origami}{center}{mincho}{center}{white}
\PatternVariant{12. ribbons（格子リボン）}{ribbons}{left}{gothic}{right}{black}

\clearpage

\begin{center}
  {\sffamily\gtfamily\bfseries\large 和柄（日本の伝統文様）}\par
  \smallskip
  {\small 線画系 6 種・面塗り系 4 種．いずれも公有の意匠を TeX で作図した独自図形}
\end{center}
\vspace{2mm}

\PatternSample{13. asanoha（麻の葉・線画）}{asanoha}{left}
\PatternSample{14. seigaiha（青海波・線画）}{seigaiha}{center}
\PatternSample{15. shippou（七宝・線画）}{shippou}{left}
\PatternSample{16. kikkou（亀甲・線画）}{kikkou}{center}

\clearpage

\PatternVariant{17. sayagata（紗綾形・線画）}{sayagata}{center}{gothic}{center}{black}
\PatternVariant{18. kagome（籠目・線画）}{kagome}{left}{mincho}{right}{black}
\PatternSample{19. yagasuri（矢絣・面塗り）}{yagasuri}{left}
\PatternSample{20. tatewaku（立涌・面塗り）}{tatewaku}{center}

\clearpage

\PatternSample{21. mitsukuzushi（三崩し・面塗り）}{mitsukuzushi}{left}
\PatternVariant{22. ajiro（網代・面塗り・白抜き題字）}{ajiro}{center}{mincho}{center}{white}

\clearpage

\begin{center}
  {\sffamily\gtfamily\bfseries\large 数学的モチーフのパターン}\par
  \smallskip
  {\small 9 patterns: truchet / isocubes / sierpinski / lattice / \ldots}
\end{center}
\vspace{2mm}

\PatternTuned{23. truchet（トゥルーシェタイル）}{truchet}
\PatternTuned{24. isocubes（等角投影の立方体）}{isocubes}
\PatternTuned{25. sierpinski（シェルピンスキー三角形）}{sierpinski}
\PatternTuned{26. lattice（格子点と基本平行四辺形）}{lattice}

\clearpage

\PatternTuned{27. moire（同心円とモアレ）}{moire}
\PatternTuned{28. envelopes（放物線の包絡線）}{envelopes}
\PatternTuned{29. girih（星形多角形の組子）}{girih}
\PatternTuned{30. packing（円充填）}{packing}

\clearpage

\PatternTuned{31. spirals（螺旋）}{spirals}

\begin{center}
  {\sffamily\gtfamily\bfseries\large 濃度・色味・寸法の制御}
\end{center}
\vspace{2mm}

\PatternTuned[pattern-tone=pale]{pattern-tone=pale}{isocubes}
\PatternTuned[pattern-tone=light]{pattern-tone=light}{isocubes}
\PatternTuned[pattern-tone=strong]{pattern-tone=strong（既定 normal は従来と同一）}{isocubes}

\clearpage

\PatternTuned[pattern-scale=1.6]{pattern-scale=1.6}{moire}
\PatternTuned[pattern-scale=0.7]{pattern-scale=0.7}{packing}
\PatternTuned[pattern-tone=light,pattern-accent=SMcoverBlue]%
  {pattern-accent=SMcoverBlue（pattern-tone=light と併用）}{girih}
\PatternTuned[pattern-tone=pale,pattern-accent=SMcoverOrange]%
  {pattern-accent=SMcoverOrange（pattern-tone=pale と併用）}{sierpinski}

\clearpage

\begin{center}
  {\sffamily\gtfamily\bfseries\large バナー以外の部品}
\end{center}
\vspace{2mm}

\MathMagazineSetup{pattern=lattice,pattern-tone=light}

\noindent
\verb|\MMPatternRule| は節の区切りに使う細い帯である．高さは省略可能引数で
指定する．

\MMPatternRule

\noindent
既定の高さは 3.2mm である．次は 6mm を指定した例．

\MMPatternRule[6mm]

\noindent
\verb|\MMPatternSwatch| は行中に置く見本である．小さな見本では
\verb|pattern-scale| を 1 未満にすると図柄が 1 単位分収まる．
\MathMagazineSetup{pattern-tone=normal,pattern-scale=0.34}%
truchet \MathMagazineSetup{pattern=truchet}\MMPatternSwatch{11mm}，
girih \MathMagazineSetup{pattern=girih}\MMPatternSwatch{11mm}，
spirals \MathMagazineSetup{pattern=spirals}\MMPatternSwatch{11mm}
のように凡例へ使える．

\end{document}
